\chapter{74}

\section{Dienstag, 4. September}



Kurz vor Mitternacht hatte auch Louis in den Schlaf gefunden. Ava war
ihm fremd. Er hatte es geschafft, dass sie bereits schlief, als er sich
hinlegte. Und bevor sie aufwachte, stand er geduscht in der Küche und
bereitete das Frühstück vor.

Aus der Praxis vernahm er ein kurzes Blöken, das sich mit Manuelas
Stimme mischte. Sie klang besänftigend. Plötzlich überraschte ihn ein überwältigendes Gefühl von Geborgenheit. Zum wiederholten Mal hier oben. Es lag an dieser Fürsorge. Daran, wie freundlich Joh und Manuela mit Menschen und auch mit ihren Tieren umgingen. Er hatte ein verdammtes Glück, die \emph{Sunnewaid} gefunden zu haben. Wie hatte er sich eben noch nach geeigneteren Fluchtplänen sehnen können?

Ava stand verschlafen im Türrahmen. 

«Hi, Ciro, guten Morgen.»

Sie wirkte lockerer als am Vorabend. 

\sectionbreak

Der restliche Tag hielt Louis stabil an der Oberfläche. Die Aufgaben
jagten sich. Gleich nach dem Frühstück folgte er Manuela in die Praxis.
Er half ihr, das Schaf auf den OP-Tisch zu heben. Sah zu, wie sie es
versorgte und staunte später, dass das Tier bereits wieder selbständig
aufstehen und herumgehen konnte. \emph{Tosca} streckte ununterbrochen
ihren schwarzen Kopf nach vorne, blickte Louis und Manuela aus
kreisrunden Knopfaugen abwechselnd an und wackelte mit den Ohren. Gleichzeitig
schwankte auch der dicke Bauch hin und her. Es sah zu lustig aus und
liess die beiden vergnügt lachen. Manuela beugte sich vor, umschlang
\emph{Toscas} wilden Lockenkopf und schloss die Augen.

«Ich glaube, wir haben es geschafft, meine Süsse.»

Dann ging Manuela Louis voran in den Stall. Ava war bereits da. 
Manuela reichte den beiden Schreibblock und Bleistift und setzte sich
inmitten des Stalls auf eine Holzkiste.

«Vielleicht wollt ihr euch die einzelnen Aufgaben notieren,» sie hielt ihren Bleistift 
in die Höhe, «bis morgen Abend sollte alles geschafft sein. Die Vorbereitung des Stalls hat Priorität. 
\emph{Tosca} wird einziehen,» das Tier hatte sich neben Manuelas Beinen zu Boden gelegt und blickte
schräg zu ihr hoch, sie streichelte ihm über den Rücken, «ein bisschen verloren vorerst, so ganz alleine, mein Liebling - weiter geht es um die Ernte im Garten, um die Verarbeitung der Früchte und Gemüse, um unsere Kocheinsätze, um Wäsche und Einkäufe und was sonst noch zu tun ist. Zu dritt werden wir die Dinge schaffen. Ich bin sehr froh, dass ihr da seid.»

Es dauerte, bis Louis und Ava den Arbeitsplan kannten und ihre Fragen
geklärt waren.
Dann war Morgenpause-Zeit. Louis bereitete Kaffee. Ava schnitt Brot und
Manuela stellte Johs Nummer ein. Sie nagte an ihrer Unterlippe und
wartete.

«Joh? - schön, deine Stimme zu hören. Ja, du, - und bei euch?»

Sie schaute konzentriert in die Ferne.

«Ja, ich bin froh, dass es so leicht ging, - hm, Ciro und Ava haben
geholfen, - doch, gut, das Kleine ist wohlauf, - genau, die Werte sind
okay, \emph{Tosca} spaziert schon wieder rum mit ihrem Babybauch,» jetzt
lachte Manuela, «doch, ich denk, sie kann es halten. Hier habe ich sie
unter Kontrolle, - genau, also, -»

Manuela hörte zu.

«Ja, finde ich gut, ja - dann bleibe ich mit Ava unten und Ciro kommt
übermorgen hoch zu euch? - gute Idee, machen wir so. Und, Joh, warte -»

Sie nickte Louis zu und er reichte ihr stumm eine Espressotasse.

«- ja, ja, - werde ich Ciro ausrichten, ciao, mein Liebster, - ciao.»

Sie beendete den Anruf und wandte sich Louis zu.

«Am Wochenende wird Joh mit dir den ganzen Abstieg zu Fuss machen, Ciro.
Er wird dir den Weg zeigen, den wir mit den Schafen auf die Bettmeralp
runter machen zur \emph{Schafscheid}, okay?»

«Ja, darauf freue ich mich.»

Louis freute sich wirklich. Johs Lebenskonzept interessierte ihn je länger je mehr. Während des gemeinsamen Wanderns würde Louis ihm bestimmt die eine oder andere Frage stellen können.

Louis reinigte den OP-Tisch in der Praxis und füllte den Seifenspender
auf. Draussen hörte er Stimmen und zwischen Bodenwischen und
Eierlieferung vibrierte sein Handy. Er wischte sich die Hände ab, 
zog das Gerät aus der Hosentasche und schaute auf das Display.

«Nein.»

Er rollte die Augen und ärgerte sich. Sie sollte hier nicht anrufen.
Er hatte es ihr doch deutlich gesagt. Louis schaute um sich. Manuela war im Garten. Im Wohnraum nebenan musste Ava am Staubsaugen sein. Er schloss die Tür, stellte sich an die Wand und
drückte auf den grünen Knopf.

«Louis? Tut mir leid, ich musste dich anrufen.»

«Valentina?» er zögerte und sprach leise, «was ist denn?»

«Von Maria weiss ich, dass Giorgia aufgewacht ist. Die Salas sind
erleichtert, nur, Giorgia soll null Erinnerung an den Vorfall haben.
Dieser Mantovani rief wieder an, Louis,» Valentina schnappte nach Luft
und redete dann weiter wie ein Wasserfall, «der Anruf hat mich so
erschreckt. Er wollte schon wieder wissen, ob ich weiss, wo du bist. Die
suchen dich immer noch als Zeugen. Wollen wissen, ob du, ob du bei
Giorgia warst, ob du jemanden gesehen hast, ich - »

«Und was hast du gesagt?»

«Noch mal das gleiche, aber, ich kann bald nicht mehr, und, und warum
ich vor allem anrufe, Louis, über die Medien werden nach wie vor Zeugen
dieser Nacht gesucht und speziell auch du.»

Louis’ Hals schnürte sich zu, als würde ein satt umgelegter Strick
ruckartig zusammengezogen.

«Was, wie?»

Valentinas Stimme zitterte.

«Mit Bild, irgend einem alten, mit Balken über deinen Augen und dein
Name steht da, mit deinen Initialen L.M.V., und mit der Frage, ob jemand
weiss, wo du dich aufhältst.»

Er versuchte, seine Erregung zu kontrollieren, klang aber dennoch
schnittig.

«Und warum werde ich gesucht?»

«Na, eben, als Zeuge. Moment, bei mir klingelt es an der Haustür.»

Über die Leitung waren nun Schritte zu hören, das Lösen einer Kette und
das Quietschen einer Türklinke. Dann eine männliche Stimme.

«Hi, Valentina, stör ich?»

«Du? Hi.»

Valentinas Stimme hörte sich gebrochen an.

«Moment, bin gerade am Telefon, komm rein - ja, da bin ich wieder. Ich
rufe dich später zurück,» ihre Stimme wurde ganz leise, «Bastiano ist
hier,» und damit war die Leitung unterbrochen. 
Louis hielt sich am Türrahmen fest. Er befürchtete, umzukippen. Was zum Geier
war das denn? Bastiano? \emph{Der} bei Valentina? Jetzt verstand er gar
nichts mehr. Die alte, bohrende Wut kam wieder hoch. Entgegen jeder
Ablehnung hing der Typ wie eine Klette an ihm. Was hatten die zusammen
zu schaffen, Valentina und der? Nun war es soweit, er wurde gesucht, als
Zeuge? Verdammt. Louis zitterte und beim Einstecken des Handys kam er versehentlich an
den Rückrufknopf. Nein. Erschrocken beendete er die Verbindung.
Von draussen hörte er Manuela «Ciro» rufen.
